\section{Twenty-seven unrestricted integers by an odd scale}
\label{sec:integer27}\label{sec:integer}

The ordinary and normalized fifteen-coordinate certificates admit a
smaller integer-domain conversion than the seven-positive-port conversion
in package~14. Require positivity only for
\[
 J,F,\alpha,z,\eta,\zeta,
\]
put $i=2r+1$, $s=2v+1$, and $w=u^2$, and allow
$r,v,u,f,h,T,\tau,y,\gamma$ to be arbitrary integers. Replacing each of
the six positive ports by its own copy of
\[
 \operatorname{Pos}(x,y,z)=1+x^2+y^2+z(z+1)
\]
gives $6\cdot3+9=27$ unrestricted integer unknowns. The map is positive
and onto the positive integers: represent $4n-3$ by three squares and
separate its one odd square from its two even squares.

The new point is the signed scale $s$. The source formulas still give
$q\ge16$, $C\ge0$, $q(q^2-1)<R<q^4-2$, and $|Y|\ge q^3$.
The usual unit arguments use squares and remain valid for signed roots.
The main norm excludes $X=0$: then $D_{\rm main}=ac+\gamma H$ and
$\Delta=a^2+H$ would imply $0\equiv1\pmod H$. Hence $X\ge q$.

Suppose $Y=-Z<0$. Define
\[
 B=(X+1)Z-2,\qquad c_0=-c=kZ-\eta,\qquad P_0=2XZ^2+1.
\]
Then $k(Z-1)<c_0<kZ$, $P_0>B$, and $c_0>k$. The first and main
Pell equations yield
\[
 k=2\psi_{P_0}(n),\qquad c_0=\psi_B(p),\qquad p>n,
 \qquad 2n=R+\varepsilon+v_0|E|,
\]
where $v_0\ge0$ and $\varepsilon\in\{-1,1\}$. Thus
$n\ge(R-1)/2\ge24$, $c_0>B(B^2-1)^2$, and $c_0>2p,2R$.
The normalized strong-rank argument, or its ordinary counterpart, and
the two signed auxiliary congruences give $p=R$. The sign change
$c=-c_0$ changes only a sign removed when the first congruence is squared.
The bound $|E|>R+2$ excludes wrapping. If $\varepsilon=-1$, Pell
duplication and $\chi_B(2)>P_0$ give
$c_0=\psi_B(2n+1)>Bk>kZ$, a contradiction. Therefore
$R=2n-1=2r_0+1$.

The elementary Pell growth estimates now imply
\[
 Z>\frac{c_0}{k}\ge
 \frac12\left(\frac{(2B-1)^2}{2P_0}\right)^{r_0}
 >\frac{X^{r_0}}2>2^R.
\]
The strict middle inequality is justified by
\[
 (2B-1)^2-2P_0X
 =4(2X+1)Z^2-20(X+1)Z+25-2X>0.
\]
At $Z=4$ its value is $46X+9$, and it increases thereafter; here
$X\ge16$ and $Z\ge4096$.
For the original negative Pell parameter $a+2=-B$, the two Pell projections
modulo $H=4a+3$ remain $2^R$ and $2^{-R}$. Since $R$ is odd and
$c=-\psi_B(R)$, the two signs of the main root respectively imply
\[
 X+2^R\equiv0\pmod H
 \quad\hbox{or}\quad
 X2^R\equiv1\pmod H.
\]
They contradict $0<X+2^R<|H|$ and $1<X2^R<XZ<|H|$. Thus $s>0$.
The inherited signed compiler theorem now applies. This exclusion only
needs a nonzero strong multiplier; the odd parametrization guarantees
that requirement.

Completeness follows from an exact parity identity, rather than an
additional padding conjecture. Write $q=2^t$, $M_0=q^2$, and
\[
 S_0=(C-1)+qF,\qquad
 T_0=(M_CJ+1)+q(M_F^{\rm nat}J-1),
\]
where $M_F^{\rm nat}$ is the unshifted mask. Native mask disjointness gives
$S_0\mathbin{\&}T_0=0$, and
$\operatorname{popcount}(T_0)=t+2$. Consequently
\[
 R=(M_0-1-S_0)M_0+(S_0+T_0),\qquad
 \operatorname{popcount}(R)=3t+2.
\]
The half-binomial identity gives $v_2(Y)=3t$, so $s=Y/q^3$ is odd.
The existing odd-multiplier and square-$w$ completeness choices are
unchanged. The new substitution costs $6\cdot7+2+2+1=47$ gates.

\begin{center}
\begin{tabular}{lrrr}
\hline
Construction & Equations & Degree & Gates\\
\hline
Normalized product & 1 & 274 & 114\\
Reduced normalized product & 1 & 216 & 115\\
Ordinary product & 1 & 172 & 116\\
Shifted ordinary system & 6 & 46 & 107\\
Ordinary sum of squares & 1 & 92 & 124\\
\hline
\end{tabular}
\end{center}
All rows use 27 unrestricted integers. The full signed proof, emitted
circuits, and reproducible arithmetic checks accompany the report in
\src{proofs/ODD_SCALE_27_PROOF.md}. The inherited signed and ordinary
rank arguments are those of packages 14 and 12~\cite{signed29,gap}.

The degrees are uniform over every fixed $\beta>0$. To see this without
a numerical coefficient specialization, write $p_*$ for the highest
homogeneous part of a substituted port, and set
\[
 Q_*=\beta J_*,\quad k_*=\eta_*+\zeta_*,\quad
 C_*=Q_*-F_*-\alpha_*,\quad
 N_{t,*}=w_*C_*-z_*Q_*.
\]
Here $w_*=u^2$, $s_*=2v$, and $i_*=2r$. Each guarded highest part is
a sum of three squares in its own disjoint variables. In particular
$N_{t,*}\ne0$, because its specialization at $u=0$ is $-z_*Q_*\ne0$.
The highest parts of the three product polynomials in table order are
\[
\begin{aligned}
 -8Q_*^{91}h\gamma i_*^4k_*^{13}w_*^{13}s_*^{26}T^2f^2N_{t,*},\\
 -8Q_*^{71}h\gamma i_*^2k_*^9w_*^{11}s_*^{20}T^2f^4N_{t,*},\\
  8Q_*^{55}h\gamma i_*^2k_*^9w_*^7s_*^{16}T^2f^4N_{t,*}.
\end{aligned}
\]
These are nonzero polynomials of degrees 274, 216, and 172. The unique
degree-46 system residual has highest part $a_*^2c_*^2T^2f^4$,
where $a_*=w_*s_*Q_*^4$ and $c_*=k_*s_*Q_*^3$; its square gives
the exact SOS degree 92. No value of the other fixed coefficient ports
affects these conclusions.

\section{Twenty-six unrestricted integers by robust separating weights}
\label{sec:integer26}

The twelve-coordinate norm also admits an integer-domain refinement.
Keep seven positive ports
$J,F,\alpha,z,i,\eta,\zeta$, replace them by disjoint copies of
$\operatorname{Pos}$, put $s=2v+1$ and $w=u^2$, and leave $h,T,t,v,u$
unrestricted. This uses $7\cdot3+5=26$ integers. In this construction
$i$ is positive, with no assertion that it is odd. That distinction is
necessary because the normalized multiplier is constrained to be
divisible by $H$.

Use the common prefix and write
\[
 \begin{gathered}
 P=X+ac,\quad S=1+\Delta c^2,\quad
 L=EkY,\quad A=1+L(L+k),\\
 G=Hic^2,\quad D=1+\Delta G^2,\quad
 Q=\Delta(D-1),\quad d=Q-1,\\
 v_0=cT,\quad b=c+RD,\quad
 e=Q(Dv_0^2+b^2)-1,\quad B_0=2Qv_0b.
 \end{gathered}
\]
Replace the positive-domain separating weights by
\[
 \begin{split}
 U_0&=(T^2+1)(a^2+5)H^2ic^4,\\
 M&=1+2(U_0+b),\qquad N=1+H^2M(A+1).
 \end{split}
\]
Use these weights in the same explicitly expanded two-quadratic-norm
polynomial, denoted $\mathcal N_{26}$. The formulations are
\[
 N_k=1,\quad N_t=1,\quad\mathcal N_{26}=0,
 \qquad\hbox{or}\qquad
 N_kN_t(1+\mathcal N_{26})-1=0.
\]

Here are the additional signed-domain details. The unit guard is valid
because the inner norm coefficient is even, so $\mathcal N_{26}$ is a
square modulo four. Weak transport therefore gives $C\ge0$ and
$0<R<q^4-2$ before roots are recovered. Even if $X=0$, the odd scale
ensures $|a|\ge q^3$, $|c|>1$, and
\[
 \Delta\ge(q^3-1)(q^3-3)>q^4>R,
 \qquad S>1,\quad A\ge1.
\]
Also $D>\Delta+|c|$, which gives the crucial interval
\[
 1<b<Q,
 \qquad Q-b=(\Delta-R)D-\Delta-c>0.
\]

Consider the algebraic triples satisfying
\[
 (P+H\gamma)^2=S,\qquad \tau^2=A,\qquad
 (dy^2-e)^2=B_0^2D.
\]
Every complex auxiliary branch obeys
$|y|<2(|v_0|\sqrt D+b)<M$. Indeed
$\sqrt D<|a+2|\,|G|$, and the factors in $U_0$ bound respectively
$|T|$, $|a+2|$, $|H|$, and $|c|^3$ by integer polynomials.
Furthermore
\[
 \frac{N\sqrt S}{|H|}>M(A+1)
 \ge M\sqrt A+M>M\sqrt A+|y|.
\]
Thus $(\gamma,\tau,y)\mapsto y+M\tau+N\gamma$ is injective on
distinct triples, including at $A=1$ and when the auxiliary quartic
has repeated roots. The exact norm factorization and Galois invariance
of the supplied integer $t$ imply that $\gamma,\tau,y$ are rational.

Both $\tau$ and $P+H\gamma$ are integers, so the denominator of
$\gamma$ divides $H$. For $y=r/s$ in lowest terms, the quartic
\[
 d^2y^4-2dey^2+e^2-B_0^2D=0
\]
shows that every prime dividing $s$ divides $d$. Because $Q$ is
divisible by $H^2$, $\gcd(d,H)=1$. The identity
$y+N\gamma=t-M\tau$ therefore makes $y$ integral, and
$N\equiv1\pmod H$ then makes $\gamma$ integral. This argument
does not assume that $T$ or $B_0$ is nonzero.

If $T=0$, then $dy^2=Qb^2-1$, hence
\[
 d(y^2-b^2)=b^2-1>0.
\]
Integrality gives $|y|\ge b+1$, whereas $b<Q$ gives $d\ge b$.
The left side is then at least $b(2b+1)>b^2-1$, a contradiction.
Therefore $B_0\ne0$, and
\[
 f=\frac{e-dy^2}{B_0},\qquad f^2=D
\]
is an integer. The full normalized strong and auxiliary equations have
now been recovered. The main norm excludes $X=0$ modulo $H$. The
negative-scale argument of the preceding section applies to the nonzero
normalized multiplier $Hi$ and excludes $s<0$. The inherited signed
compiler theorem then proves soundness.

For completeness, retain the canonical square $w$ and odd $s$. Choose
the strong Pell index to be a common multiple of $R$ and a period of
the invertible Pell matrix modulo $Hc^2$. This gives a positive
$i=\psi_{a+2}(m)/(Hc^2)$ and the required auxiliary quotient, without
requiring odd $i$. The canonical roots make
$t=y+M\tau+N\gamma$ an integer, and all seven positive ports have
$\operatorname{Pos}$ representations.

The exact integer degrees are
\[
 \deg d=106,\quad \deg b=\deg M=92,\quad\deg N=154,
 \quad \deg(P^2-S)=29.
\]
The two main-root branches have degrees $9$ and $-2$; therefore the
$N\gamma$ term dominates every linear factor of the norm. Its leading
form is
\[
 (\mathcal N_{26})_{3544}
 =d_{\rm top}^{8}N_{\rm top}^{16}
   (H_{\rm top}c_{\rm top}^{2})^8\ne0.
\]
The comparison system has degree 3544 and 195 gates; the product has
degree 3560 and 199 gates. A sum of squares has degree 7088 and 202
gates. All counts include the 52 substitution gates. These constructions
provide a variable-count improvement at a substantial degree and circuit
cost; the 27-integer degree-92 polynomial remains a useful alternative.
